\begin{proof}[Proof of \cref{obj:lem:cty-support-boundary-angular-packing}]
Fix \(q\ge1\) and \(L\ge4\), and write \(S=[-1/2,1/2]^2\).

\emph{Step 1: the fixed profiles.}  Let \(\varphi:\mathbb R^2\to[0,1]\) be the normalized radial bump, that is, the smooth radially symmetric function with
\[
\varphi(0)=1,
\qquad
0\le\varphi\le1,
\qquad
\varphi(z)=0\ \text{ whenever }\ \lVert z\rVert_2\ge1,
\]
and let \(\Psi:\mathbb R\to[0,1]\) be the smooth transition function, which vanishes on \((-\infty,0]\) and equals one on \([1,\infty)\).  For each integer \(j\ge0\), fix a nonnegative derivative envelope \(b_j\) satisfying
\[
\lVert D^j\varphi(z)\rVert\le b_j
\qquad\text{for every }z\in\mathbb R^2 .
\]
These envelopes are finite fixed constants attached to the bump profile.
% lean: packingBump, packingBump_zero, packingBump_mem_Icc, packingBump_eq_zero_of_one_le_norm, packingBumpDerivativeBound, packingBump_iteratedFDeriv_le_derivativeBound

\emph{Step 2: the constants.}  Define the smoothness-dependent derivative scale and the one-observation divergence constant by
\[
D_q=1+q+1310720+\sum_{j=0}^{q}b_j ,
\qquad
C_{\mathrm{KL}}=1310720 ,
\]
and set
\[
c_0=\frac1{24},\qquad
c_1=\frac1{1024D_q},\qquad
c_w^-=c_w^+=1,\qquad
c_{\mathrm{rad}}=\frac18,\qquad
\alpha=\frac1{16}.
\]
Since every \(b_j\ge0\), the scale \(D_q\) is finite and satisfies
\[
D_q\ge1,
\qquad
D_q\ge q,
\qquad
D_q\ge C_{\mathrm{KL}} .
\]
Hence all six constants are positive and \(\alpha=1/16<1/8\).
% lean: packingBumpDerivativeScale, packingBumpDerivativeScale_pos, angularPackingOnePointKLConstant, packingBumpDerivativeBound_le_scale

The three displayed inequalities give the divergence budget used below:
\[
256\,q\,C_{\mathrm{KL}}
\le
256\,D_q^{2}
\le
(1024D_q)^{4},
\]
the first step because \(q\le D_q\) and \(C_{\mathrm{KL}}\le D_q\), the second because \(D_q\ge1\).
% lean: cty_support_boundary_angular_packing

\emph{Step 3: the eventual regime.}  With \(a_n=(\log n/n)^{1/4}\) as in \cref{obj:def:frontier-rates}, set
\[
w_n=a_n^{1/q},
\qquad
\delta_n=\frac{a_n}{1024D_q},
\qquad
M_n=\left\lfloor\frac1{12w_n}\right\rfloor .
\]
% lean: angularGridRadius, angularPackingScaledDelta, angularGridSize
Because \(a_n\to0\), for all sufficiently large \(n\) we have \(n\ge2\), \(a_n\le1\), and
\[
0<w_n\le\frac1{24},
\qquad
0<\delta_n\le\frac18 .
\]
% lean: angularGridRadius_eventually_small, angularPackingScaledDelta_eventually_admissible
Moreover, since \(n a_n^{4}=\log n\),
\[
n\,C_{\mathrm{KL}}\,\delta_n^{4}
=
C_{\mathrm{KL}}\,\frac{\log n}{(1024D_q)^{4}}
\le
\frac{\log n}{256\,q}
\qquad\text{by the budget of Step 2,}
\]
while \(M_n\ge 1/(24w_n)=\tfrac1{24}a_n^{-1/q}\) gives \(\log M_n\ge \tfrac1{4q}\log n-o(\log n)\); hence, enlarging \(n\) once more,
\[
n\,C_{\mathrm{KL}}\,\delta_n^{4}\le\alpha\log M_n .
\]
% lean: angularPackingScaledDelta_fourth_power, angularPackingScaledDelta_eventually_klBudget
Finally, for every \(r\ge c_{\mathrm{rad}}a_n=a_n/8\),
\[
2\cdot(8\delta_n)\le\frac{r}{8},
\]
because \(D_q\ge1\) makes \(16\delta_n=a_n/(64D_q)\le a_n/64\le r/8\).
% lean: angularPackingScaledDelta_cutoff_fully_active

Fix any \(n\) in this eventual regime; all remaining claims are verified at that \(n\).

\emph{Step 4: the packing points and cells.}  Take the equispaced grid on the middle half of the lower edge of \(S\),
\[
x_j=\left(-\frac14+\frac{j}{2(M_n+1)},\,-\frac12\right),
\qquad j=1,\ldots,M_n,
\]
and let \(C_j=B_2(x_j,w_n)\cap S\) as in the statement, with \(B_2\) the closed Euclidean ball of \cref{obj:synth_43}.  Each \(x_j\) has second coordinate \(-1/2\) and first coordinate of modulus strictly below \(1/4\), so \(x_j\in\operatorname{bd}(S)\); and since \(w_n\le1/24\), each \(C_j\) is the closed upper half-disc \(\{z\in B_2(x_j,w_n):z^{(2)}\ge-1/2\}\), where \(z=(z^{(1)},z^{(2)})\).
% lean: angularGridCenter, angularGridCenter_mem_packingSquare_frontier

From the definition of \(M_n\), \(M_n\le 1/(12w_n)\).  The small-radius bound \(w_n\le1/24\) also gives \(1\le1/(12w_n)\), and therefore
\[
M_n+1\le \frac{1}{12w_n}+\frac{1}{12w_n}=\frac1{6w_n}.
\]
Thus the grid spacing obeys
\[
3w_n\le\frac1{2(M_n+1)} ,
\]
so that for \(i\ne j\)
\[
\lVert x_i-x_j\rVert_2\ge 3w_n\ge w_n=c_w^-a_n^{1/q},
\]
and the closed cells \(C_i,C_j\), of radius \(w_n\) around centers at distance at least \(3w_n\), are disjoint.
% lean: angularGridSize_spacing, angularGridCenter_three_radius_separated, angularGridPackingCells_disjoint
The radius bounds \(c_w^-a_n^{1/q}\le w_n\le c_w^+a_n^{1/q}\) hold with equality since \(c_w^\pm=1\).  For the packing size, \(w_n\le1/24\) forces \(1/(12w_n)\ge2\), and \(\lfloor u\rfloor\ge u/2\) for \(u\ge2\), so
\[
M_n\ge\frac1{24w_n}=\frac1{24}\,a_n^{-1/q}=c_0a_n^{-1/q}.
\]
% lean: angularGridSize_lower, angularGridSize_frontier_lower

\emph{Step 5: the laws \(P_\omega\).}  Fix \(\omega\in\{0,1\}^{M_n}\).  Define the regression profile on \(\mathbb R^2\) by
\[
\mu_\omega(x)
=
\frac12+\frac{x^{(1)}}{8}
+
\sum_{j=1}^{M_n}\omega_j\,\delta_n\,
\varphi\!\left(\frac{x-x_j}{w_n}\right),
\]
% lean: packingAffineBaseline, localizedPackingBump, packingRegression
and the design density on \(S\) by
\[
f_\omega(x)
=
1+\sum_{j=1}^{M_n}\omega_j\,
t_n\bigl(\lVert x-x_j\rVert_2\bigr)\,\kappa_j(x),
\qquad
\kappa_j(x)=
\begin{cases}
\dfrac{x^{(1)}-x_j^{(1)}}{\lVert x-x_j\rVert_2}, & x\ne x_j,\\[2mm]
0,& x=x_j,
\end{cases}
\]
where the radial tilt is
\[
t_n(r)
=
-\,\frac{2\delta_n\,\varphi\!\left(\tfrac{r}{w_n}e_\circ\right)\,\chi_n(r)}
{\max\{r/8,\ 8\delta_n\}},
\qquad
\chi_n(r)=\Psi\!\left(\frac{r}{64\delta_n}-1\right),
\]
and \(e_\circ\) is a unit vector of \(\mathbb R^2\) (the bump is radial, so the choice is immaterial).
% lean: angularCutoff, angularRadialProfile, angularTilt, packingDirectionCos, packingAngularTerm, packingAngularDensity, angularDesignDensity
Let \(P_\omega\) be the joint law on observation coordinates obtained as follows.  Its \(X\)-marginal is the measure with Lebesgue density \(f_\omega\mathbf 1_S\), and at score \(x\) the outcome kernel is the probability measure
\[
\mu_\omega(x)\,N(1,1)+\{1-\mu_\omega(x)\}\,N(0,1).
\]
Equivalently, \(P_\omega\) is the image under \((x,y)\mapsto(y,x)\) of the product-kernel measure generated by this \(X\)-marginal and this outcome kernel.  The declared regression and variance profiles of the resulting law are
\[
\mu_{P_\omega}(x)=\mu_\omega(x),
\qquad
\sigma^2_{P_\omega}(x)=1+\mu_\omega(x)\{1-\mu_\omega(x)\}.
\]
They are \(X\)-marginal-almost-everywhere versions of the conditional mean and variance in the sense of \cref{obj:def:cty-nonparametric-class}.
% lean: bernoulliGaussianLaw, bernoulliGaussianKernel, jointBernoulliGaussianLaw, bernoulliGaussianCtyLaw
Since \(|x^{(1)}|\le1/2\) on \(S\), the bumps have disjoint supports \(B_2(x_j,w_n)\), and \(\delta_n\le1/8\),
\[
\frac14\le\frac12-\frac1{16}\le\mu_\omega(x)\le\frac12+\frac1{16}+\delta_n\le\frac34
\qquad\text{for }x\in S,
\]
so the mixture weight is legitimate and the clipping of \(\mu_\omega\) to \([0,1]\) needed off \(S\), and to \([1/4,3/4]\) used in Step 10, are both silent on \(S\).
% lean: clippedPackingRegression, clippedPackingRegression_eq_on_square, middleHalfClip_clippedPackingRegression_ae_eq

\emph{Step 6: model class, support, and cell masses.}  The angular tilt obeys \(|t_n(r)\kappa_j(x)|\le1/4\) uniformly, and the bumps attached to distinct centers have disjoint supports, so
\[
\frac34\le f_\omega(x)\le\frac54
\qquad (x\in S),
\]
whence \(L^{-1}\le f_\omega\le L\) for \(L\ge4\); \(f_\omega\) is continuous on \(S\), and \(S\) is compact with \(S\subseteq[-L,L]^2\) and with \(\operatorname{bd}(S)\) Lipschitz-parametrized by \([0,1]\).  The conditional variance \(1+\mu_\omega(1-\mu_\omega)\) is continuous and lies in \([1,5/4]\subseteq[L^{-1},L]\).  Finally, for \(0\le j\le q\),
\[
\delta_n\,w_n^{-j}\,b_j
=
\bigl(a_nw_n^{-j}\bigr)\cdot\frac{b_j}{1024D_q}
=
w_n^{\,q-j}\cdot\frac{b_j}{1024D_q}
\le 1,
\]
using \(w_n^{q}=a_n\), \(w_n\le1\) and \(b_j\le D_q\); since the \(j\)th derivative of the localized bump \(\delta_n\varphi((\cdot-x_j)/w_n)\) is bounded by \(\delta_nw_n^{-j}b_j\), these bounds give \(\lVert D^j\mu_\omega\rVert\le L\) on \(S\) for \(0\le j\le q-1\) together with \(L\)-Lipschitz continuity of \(D^{q-1}\mu_\omega\), that is
\[
\mu_\omega\in\mathcal H^q(S,L)
\]
in the sense of \cref{obj:synth_36}.  Consequently \(P_\omega\in\mathcal P_{\mathrm{NP}}(L,q)\) as in \cref{obj:def:cty-nonparametric-class}, with \(\mathcal X_{P_\omega}=S\).
% lean: angularDesignDensity_mem_Icc, angularPackingScaledDelta_derivative_scaling, clippedPackingRegression_mem_holder_of_scaling_bounds, angularPackingCtyLaw_mem_class_scaledDelta, angularPackingCtyLaw_support

For the cell masses, the only \(\omega\)-dependent part of \(f_\omega\) on \(C_j\) is \(\omega_jt_n(\lVert x-x_j\rVert_2)\kappa_j(x)\), and reflecting \(C_j\) in the vertical line through \(x_j\) changes the sign of \(\kappa_j\) while fixing \(\lVert x-x_j\rVert_2\); hence that term integrates to zero over \(C_j\) and
\[
P_\omega(X\in C_j)=\operatorname{Leb}(C_j)=m_n ,
\]
where \(m_n=\operatorname{Leb}(C_1)\) is common to all \(j\) because the cells are translates of each other along the lower edge.
% lean: angularPackingCtyLaw_cell_mass_eq_volume, angularGrid_packingCell_volume_eq

\emph{Step 7: locality in one bit.}  The functions \(\varphi((\cdot-x_i)/w_n)\) and \(t_n(\lVert\cdot-x_i\rVert_2)\kappa_i(\cdot)\) vanish outside \(B_2(x_i,w_n)\), and the cells are pairwise disjoint.  Therefore, on \(C_j\) both \(f_\omega\) and \(\mu_\omega\) depend on \(\omega\) only through \(\omega_j\), so
\[
\omega_j=\omega'_j
\quad\Longrightarrow\quad
P_\omega\big|_{\{X\in C_j\}}=P_{\omega'}\big|_{\{X\in C_j\}},
\]
and on \(S\setminus\bigcup_jC_j\) we have \(f_\omega\equiv1\) and \(\mu_\omega(x)\equiv\tfrac12+\tfrac{x^{(1)}}{8}\), so
\[
P_\omega\big|_{\{X\in S\setminus\bigcup_jC_j\}}
=
P_{\omega'}\big|_{\{X\in S\setminus\bigcup_jC_j\}}
\qquad\text{for all }\omega,\omega' .
\]
% lean: angularPackingCtyLaw_locality_certificate

\emph{Step 8: boundary signal.}  Put
\[
u_{j,0}=\frac12+\frac{x_j^{(1)}}{8},
\qquad
u_{j,1}=u_{j,0}+\delta_n .
\]
At \(x_j\) every bump attached to another center vanishes, because \(\lVert x_j-x_i\rVert_2\ge3w_n>w_n\) for \(i\ne j\), while \(\varphi(0)=1\); hence
\[
\mu_{P_\omega}(x_j)=u_{j,\omega_j},
\qquad
|u_{j,1}-u_{j,0}|=\delta_n=c_1a_n .
\]
% lean: packingRegression_center, packingCenterValue, packingRegression_eq_packingCenterValue, angularPackingCtyLaw_mu_center, packingCenterValue_scaledDelta_separation

\emph{Step 9: radial invariance.}  Let \(\omega^{(j)}\) be \(\omega\) with coordinate \(j\) flipped.  For every measurable \(A\subseteq[0,\infty)\), the reflection of Step 6 applied to the annular slice \(\{x:\lVert x-x_j\rVert_2\in A\}\) again cancels the \(\omega_j\)-dependent density term, so
\[
\mathcal L_{P_\omega}\bigl(\lVert X-x_j\rVert_2\bigr)
=
\mathcal L_{P_{\omega^{(j)}}}\bigl(\lVert X-x_j\rVert_2\bigr).
\]
% lean: angularDesignMeasure_map_distance_eq, angularPackingCtyLaw_flip_map_distance_eq

For the tail statement, fix \(r\ge c_{\mathrm{rad}}a_n\).  The last display of Step 3 gives \(16\delta_n\le r/8\), so \(r/(64\delta_n)-1\ge1\) and \(r/8\ge8\delta_n\); therefore the cutoff is fully active,
\[
\chi_n(r)=1,
\qquad
t_n(r)=-\,\frac{16\,\delta_n\,\varphi\!\left(\tfrac{r}{w_n}e_\circ\right)}{r} .
\]
% lean: angularPackingScaledDelta_cutoff_fully_active, angularTilt_eq_formula
Let \(A\subseteq[c_{\mathrm{rad}}a_n,\infty)\) be measurable and let
\[
D_{j,A}=\left\{x\in B_2(x_j,w_n):x^{(2)}>-\tfrac12,\ \lVert x-x_j\rVert_2\in A\right\}
\]
be the part of the cell \(C_j\) on which the two vertices differ.  Assume \(\omega_j=1\), so that \(\omega^{(j)}_j=0\); the reverse case is symmetric.  On \(D_{j,A}\) the difference of the two outcome-weighted densities is
\[
\mu_\omega f_\omega-\mu_{\omega^{(j)}}f_{\omega^{(j)}}
=
\delta_n\varphi\!\left(\tfrac{r}{w_n}e_\circ\right)
+
\left\{\frac12+\frac{x^{(1)}}{8}+\delta_n\varphi\!\left(\tfrac{r}{w_n}e_\circ\right)\right\}
t_n(r)\,\kappa_j(x),
\qquad r=\lVert x-x_j\rVert_2 ,
\]
because on the cell all other bumps and angular corrections vanish.
% lean: packingRegression_mul_density_flip_cell_identity, packingRegression_mul_density_flip_cell_radial_identity
Splitting the brace into the radius-only part \(\tfrac12+\tfrac{x_j^{(1)}}{8}+\delta_n\varphi(\tfrac{r}{w_n}e_\circ)\) and the remainder \(\tfrac18(x^{(1)}-x_j^{(1)})\), the first part contributes \(\int_{D_{j,A}}g(r)\kappa_j=0\) by the same horizontal reflection, and the second contributes
\[
\int_{D_{j,A}}\frac{x^{(1)}-x_j^{(1)}}{8}\,t_n(r)\,\kappa_j(x)\,dx
=
\int_{D_{j,A}}\frac{r}{8}\,t_n(r)\,\kappa_j(x)^2\,dx
=
-\int_{D_{j,A}}2\delta_n\varphi\!\left(\tfrac{r}{w_n}e_\circ\right)\kappa_j(x)^2\,dx .
\]
% lean: angularGridCenter_openRadialSet_weighted_direction_cancellation
Since the average of \(\kappa_j^2\) over each half-circle \(\{\lVert x-x_j\rVert_2=r,\ x^{(2)}>-1/2\}\) equals \(1/2\), the two remaining terms cancel:
\[
\int_{D_{j,A}}\bigl(\mu_\omega f_\omega-\mu_{\omega^{(j)}}f_{\omega^{(j)}}\bigr)\,dx=0 .
\]
% lean: angularGridCenter_radialSet_angular_outcome_cancellation, packingRegression_density_flip_integral_eq_zero_of_active
Together with the equality of radial slice masses this says that, for every measurable \(A\subseteq[c_{\mathrm{rad}}a_n,\infty)\), the two experiments assign the same mass and the same conditional-mean mass to the slice \(\{\lVert X-x_j\rVert_2\in A\}\); since the outcome mixture is determined by its conditional mean, the restricted one-point distance laws agree:
\[
\mathcal L_{P_\omega}\bigl(Y,\lVert X-x_j\rVert_2\bigr)\Big|_{\{z_2\ge c_{\mathrm{rad}}a_n\}}
=
\mathcal L_{P_{\omega^{(j)}}}\bigl(Y,\lVert X-x_j\rVert_2\bigr)\Big|_{\{z_2\ge c_{\mathrm{rad}}a_n\}} .
\]
% lean: angularPacking_flip_setIntegral_eq_of_active, angularPackingCtyLaw_flip_restrict_Ici_eq_of_active

\emph{Step 10: the compressed-distance divergence.}  Below the fully active radius the same cancellation is replaced by a localized setwise bound.  For every measurable \(A\subseteq[0,\infty)\),
\[
\left|
\int_{\{\lVert X-x_j\rVert_2\in A\}}\mu_\omega\,dP_\omega^X
-
\int_{\{\lVert X-x_j\rVert_2\in A\}}\mu_{\omega^{(j)}}\,dP_{\omega^{(j)}}^X
\right|
\le
\delta_n\,P_\omega\bigl(\lVert X-x_j\rVert_2\in A\cap[0,128\delta_n)\bigr),
\]
where \(P^X\) denotes the covariate marginal: outside the radii \(r<2\cdot(8\delta_n)/(1/8)=128\delta_n\) the cutoff is fully active and Step 9 gives exact cancellation, while on those radii the uncancelled increment is bounded by the bump amplitude per unit radial mass.
% lean: angularRadialSuccessIncrement_identity, angularRadialSuccessIncrement_abs_le, angularPacking_flip_success_setIntegral_abs_le
Moreover, since \(f_\omega\le5/4\) and a disc of radius \(R\) has area \(\pi R^2\le4R^2\),
\[
P_\omega\bigl(\lVert X-x_j\rVert_2<R\bigr)\le 5R^{2}
\qquad (R\ge0).
\]
% lean: angularDesignMeasure_radial_Iio_le

The complete radial laws of Step 9 provide a common marginal law for \(R=\lVert X-x_j\rVert_2\) under the two adjacent vertices.  Disintegrating the one-observation laws with respect to this common radial marginal gives measurable radial success parameters, clipped to \([1/4,3/4]\) and equal almost everywhere to the corresponding conditional Bernoulli weights, such that the preceding setwise bound implies
\[
|g_\omega(r)-g_{\omega^{(j)}}(r)|
\le
2\delta_n\,\mathbf 1\{r<128\delta_n\}
\quad\text{for the common radial marginal-almost every }r .
\]
Here the factor \(2\delta_n\) is the envelope used for the almost-everywhere Radon--Nikodym comparison, and the middle-half clipping is silent by Step 5.  For two Bernoulli-plus-Gaussian mixtures with success parameters in \([1/4,3/4]\), the divergence at a fixed radius is at most four times the squared difference of the success parameters.  Averaging over the common radial marginal therefore gives
\[
\operatorname{KL}\!\left(
\mathcal L_{P_\omega}\bigl(Y,\lVert X-x_j\rVert_2\bigr),
\mathcal L_{P_{\omega^{(j)}}}\bigl(Y,\lVert X-x_j\rVert_2\bigr)\right)
\le
4(2\delta_n)^{2}\,P_\omega\bigl(\lVert X-x_j\rVert_2<128\delta_n\bigr)
\le
4(2\delta_n)^{2}\cdot5(128\delta_n)^{2}
=
C_{\mathrm{KL}}\,\delta_n^{4},
\]
using \(4\cdot4\cdot5\cdot128^{2}=1310720=C_{\mathrm{KL}}\).
% lean: radialOutcomeLaw_klDiv_le_of_localized_success_bound, angularPacking_radialOutcome_klDiv_le

The distance data \(V_{n,P}(x_j)\) of \cref{obj:def:cty-distance-data} are \(n\) independent copies of the one-observation pair \((Y,\lVert X-x_j\rVert_2)\), so the divergence tensorizes; combining with the eventual budget of Step 3,
\[
\operatorname{KL}\!\left(
\mathcal L_{P_\omega}\{V_{n,P_\omega}(x_j)\},
\mathcal L_{P_{\omega^{(j)}}}\{V_{n,P_{\omega^{(j)}}}(x_j)\}
\right)
\le
n\,C_{\mathrm{KL}}\,\delta_n^{4}
\le
\alpha\log M_n .
\]
% lean: compressedSampleLaw_klDiv_le_of_onePoint_finite_bound

Steps 4--10 supply every displayed clause with the constants of Step 2, for all sufficiently large \(n\).
% lean: angularPackingAt_of_certificates, angularPackingAt_scaledDelta_of_radial_certificates, cty_support_boundary_angular_packing
\end{proof}
